\documentclass[12pt]{article}
\usepackage[utf8]{inputenc}
\usepackage[T1]{fontenc}
\usepackage[margin=1in]{geometry}
\usepackage{setspace}
\usepackage{xcolor}
\usepackage[colorlinks=true, urlcolor=blue!60!black, linkcolor=blue!60!black]{hyperref}
\setlength{\parindent}{2em}
\setlength{\parskip}{1.25em}
\renewcommand{\baselinestretch}{1.0}
\begin{document}
\begin{center}
{\LARGE \textit{City of Cleburne v. Cleburne Living Center, Inc.}, 473 U.S. 432 (1985)\par}
\vspace{0.8em}
{\large Case Brief\par}
\end{center}
\vspace{1em}

\section{Memory Jogger}

Intellectual disability is not a quasi-suspect classification triggering heightened Equal Protection scrutiny, yet a city violates rational basis review when it denies a special use permit for a group home based on irrational prejudice and unsubstantiated neighborhood fears (\href{https://www.courtlistener.com/opinion/111507/city-of-cleburne-v-cleburne-living-center-inc/}{CourtListener}; \href{https://supreme.justia.com/cases/federal/us/473/432/}{Justia}).

\section{Detailed Case Facts}

In July 1980, Jan Hannah, vice president and part owner of Cleburne Living Center, Inc.\ (CLC, later known as Community Living Concepts, Inc.), purchased a four-bedroom, two-bath house (with a half bath to be added) at 201 Featherston Street in Cleburne, Texas, intending to lease it to CLC for operation as a group home for thirteen men and women with intellectual disabilities (referred to in the 1985 opinion as ``the mentally retarded''). The home was to operate as a private Level I Intermediate Care Facility for the Mentally Retarded (ICF-MR) under a joint federal-state reimbursement program, subject to constant staff supervision and extensive federal and state regulations, including federal square-footage-per-resident rules (42 C.F.R.\ \S~442.447) (\href{https://supreme.justia.com/cases/federal/us/473/432/}{473 U.S. at 435--436, 450 \& nn.1--2}).

The Featherston Street property was located in an ``R-3'' Apartment House District. Under Section~8 of the Cleburne zoning ordinance, R-3 districts permitted by right apartment houses, multiple dwellings, boarding and lodging houses, fraternity and sorority houses, dormitories, apartment hotels, private clubs, and ``[h]ospitals, sanitariums, nursing homes or homes for convalescents or aged, other than for the insane or feeble-minded or alcoholics or drug addicts.'' Under Section~16, a special use permit---renewable annually and requiring the applicant to obtain the signatures of property owners within 200 feet---was required for ``[h]ospitals for the insane or feeble-minded, or alcoholic [sic] or drug addicts, or penal or correctional institutions'' (\href{https://supreme.justia.com/cases/federal/us/473/432/}{473 U.S. at 436 \& n.3}).

The City of Cleburne classified CLC's proposed group home as a ``hospital for the feebleminded'' and informed CLC that a special use permit was required. After the city's Planning and Zoning Commission held a hearing and voted to deny the permit, the Cleburne City Council held a public hearing and voted 3 to 1 to deny CLC's application (\href{https://supreme.justia.com/cases/federal/us/473/432/}{473 U.S. at 436--437 \& n.4}). The Council's denial rested on the negative attitudes and fears of nearby property owners and elderly neighbors, concern that students at a junior high school across the street (which itself enrolled about 30 students with intellectual disabilities) might harass the residents, the property's location on a 500-year flood plain, doubts regarding legal responsibility for residents' actions, and objections to the size of the home and number of occupants (\href{https://supreme.justia.com/cases/federal/us/473/432/}{473 U.S. at 448--450}).

The dispute arose amid the national deinstitutionalization movement of the 1970s and 1980s. Congress had enacted major federal protections---including Section~504 of the Rehabilitation Act of 1973 (29 U.S.C.\ \S~794), the Developmental Disabilities Assistance and Bill of Rights Act (42 U.S.C.\ \S\S~6010(1)--(2)), and the Education of the Handicapped Act (20 U.S.C.\ \S~1412(5)(B))---and Texas had passed the Mentally Retarded Persons Act of 1977 recognizing a right to live in the ``least restrictive setting,'' including group homes (\href{https://supreme.justia.com/cases/federal/us/473/432/}{473 U.S. at 443--444}). At the same time, local communities frequently invoked mid-century zoning ordinances containing archaic eugenic-era terminology (Cleburne's 1965 ordinance traced directly to a 1947 ordinance and a 1929 Dallas ordinance) to block community-based group homes (\href{https://supreme.justia.com/cases/federal/us/473/432/}{473 U.S. at 464 n.17}). The governing constitutional provision was the Equal Protection Clause of the Fourteenth Amendment.

\section{Procedural History}

CLC, Hannah, and Advocacy, Inc.\ (a nonprofit legal services organization for developmentally disabled persons) sued the City of Cleburne and several municipal officials in the United States District Court for the Northern District of Texas, alleging that the zoning ordinance violated the Equal Protection Clause both on its face and as applied. (CLC voluntarily dismissed a pendent claim under Texas's Mentally Retarded Persons Act after the District Court indicated abstention might be appropriate.) Although the District Court found as a matter of fact that the Featherston Street home would have been permitted by right if its residents were not intellectually disabled and that the Council's vote was motivated primarily by the residents' disability, the court held that intellectual disability was neither a suspect nor a quasi-suspect class and upheld the ordinance facially and as applied under minimum rational basis review (\href{https://supreme.justia.com/cases/federal/us/473/432/}{473 U.S. at 437 \& n.5}).

The United States Court of Appeals for the Fifth Circuit reversed, 726 F.2d 191 (5th Cir.\ 1984). Holding that intellectual disability is a quasi-suspect classification based on a history of ``unfair and often grotesque mistreatment,'' immutability, political powerlessness, and the importance of community housing, the panel applied intermediate scrutiny and invalidated the ordinance both on its face and as applied. The Fifth Circuit denied rehearing en banc over the dissent of six judges. The Supreme Court granted certiorari (469 U.S. 1016 (1984)), heard oral argument on March 18, 1985, ordered reargument on April 23, 1985, and decided the case on July 1, 1985. The Court affirmed the judgment insofar as it invalidated the ordinance as applied to the Featherston home, vacated the judgment in all other respects (including facial invalidation and quasi-suspect status), and remanded (\href{https://supreme.justia.com/cases/federal/us/473/432/}{473 U.S. at 437--439, 450}).

\section{Judicial Votes}

\begin{itemize}
  \item \textbf{Majority (6--3 on standard of review and as-applied remedy; 9--0 on judgment invalidating the permit denial as applied):} Justice Byron R. White delivered the opinion of the Court, joined by Chief Justice Warren E. Burger and Justices Lewis F. Powell, Jr., William H. Rehnquist, John Paul Stevens, and Sandra Day O'Connor.
  \item \textbf{Concurrence:} Justice John Paul Stevens filed a concurring opinion, joined by Chief Justice Warren E. Burger.
  \item \textbf{Concurrence in the Judgment in Part and Dissent in Part:} Justice Thurgood Marshall filed an opinion concurring in the judgment in part and dissenting in part, joined by Justices William J. Brennan, Jr., and Harry A. Blackmun.
\end{itemize}

\section{Holding}

\textbf{Judgment:} Affirmed in part, vacated in part, and remanded. The Fifth Circuit's judgment is affirmed to the extent it held Cleburne's zoning ordinance unconstitutional as applied to CLC's proposed group home, and otherwise vacated (\href{https://supreme.justia.com/cases/federal/us/473/432/}{473 U.S. at 450}).

\textbf{Rule Announced:} Classifications based on intellectual disability are not quasi-suspect and receive rational basis review rather than heightened Equal Protection scrutiny. Under rational basis review, however, a governmental classification must be rationally related to a legitimate state interest; ``mere negative attitudes, or fear, unsubstantiated by factors which are properly cognizable in a zoning proceeding,'' and ``irrational prejudice'' against intellectually disabled persons are not legitimate governmental interests (\href{https://supreme.justia.com/cases/federal/us/473/432/}{473 U.S. at 442, 446--450}).

\textbf{Holding vs.\ Dicta:} The Court's rejection of quasi-suspect status for intellectual disability and its invalidation of the special use permit requirement as applied to 201 Featherston Street are central holdings. Justice Marshall argued in dissent that because the Court struck down the permit denial under rational basis review, its lengthy Part~III rejection of heightened scrutiny was unnecessary ``two for the price of one'' constitutional adjudication akin to dicta (\href{https://supreme.justia.com/cases/federal/us/473/432/}{473 U.S. at 456--458}), though subsequent Supreme Court precedent (such as \textit{Board of Trustees of the University of Alabama v. Garrett}) treated \textit{Cleburne}'s rational basis rule for disability classifications as binding holding. The Court's observations regarding other potentially non-suspect groups (such as ``the aging, the disabled, the mentally ill, and the infirm'') and its reservation of whether the ordinance might be facially valid in other applications are dicta (\href{https://supreme.justia.com/cases/federal/us/473/432/}{473 U.S. at 445--447}).

\section{Analysis of Opinions}

\subsection{Majority Opinion (Justice White)}

Justice White began by summarizing the Court's tiered Equal Protection framework: social and economic legislation is presumed valid if rationally related to a legitimate state interest, whereas classifications based on race, alienage, or national origin trigger strict scrutiny, and classifications based on sex or illegitimacy trigger intermediate scrutiny (\href{https://supreme.justia.com/cases/federal/us/473/432/}{473 U.S. at 439--441}). Drawing on \textit{Massachusetts Board of Retirement v. Murgia}, 427 U.S. 307 (1976), White offered four reasons why intellectual disability is not a quasi-suspect class: (1) individuals with intellectual disabilities have ``a reduced ability to cope with and function in the everyday world'' across a wide spectrum of abilities, making their differences relevant to legitimate legislative goals and better suited to professional-guided policymaking than judicial oversight; (2) extensive federal and state remedial legislation demonstrates that lawmakers are actively addressing their needs rather than acting out of continuing antipathy, and heightened scrutiny might deter beneficial tailored programs; (3) this legislative responsiveness negates any claim of political powerlessness; and (4) granting quasi-suspect status would leave courts with no principled basis to deny heightened scrutiny to other large, amorphous groups such as ``the aging, the disabled, the mentally ill, and the infirm'' (\href{https://supreme.justia.com/cases/federal/us/473/432/}{473 U.S. at 442--446}).

Turning to the ordinance as applied, however, Justice White emphasized that rational basis review still forbids classifications whose relationship to an asserted goal is so attenuated as to be arbitrary or irrational, or that rest on ``a bare . . . desire to harm a politically unpopular group'' (\href{https://supreme.justia.com/cases/federal/us/473/432/}{473 U.S. at 446--447}, quoting \textit{USDA v. Moreno}, 413 U.S. 528, 534 (1973)). Because Cleburne allowed apartment houses, boarding houses, fraternity houses, dormitories, hospitals, and nursing homes in the R-3 zone without a special permit, the city had to show how the Featherston home posed a special threat not posed by those permitted uses. White systematically dismantled each of the city's justifications: neighbor hostility and fear cannot be given legal effect (\textit{Palmore v. Sidoti}); fear of harassment by junior high students was contradicted by the school's own 30 intellectually disabled students; the 500-year flood plain, legal-liability, density, congestion, and fire-hazard concerns applied equally to permitted nursing homes, hospitals, and fraternity houses. Concluding that the permit requirement rested on ``an irrational prejudice against the mentally retarded,'' the Court invalidated it as applied without reaching facial validity (\href{https://supreme.justia.com/cases/federal/us/473/432/}{473 U.S. at 447--450}).

\subsection{Concurring Opinion (Justice Stevens)}

Justice Stevens, joined by Chief Justice Burger, joined the Court's opinion but wrote separately to reject the rigid three-tier Equal Protection framework altogether. In Stevens's view, the Court's precedents actually reflect ``a continuum of judgmental responses'' governed by a single overarching rationality standard asking whether ``an impartial lawmaker could logically believe that the classification would serve a legitimate public purpose that transcends the harm to the members of the disadvantaged class'' (\href{https://supreme.justia.com/cases/federal/us/473/432/}{473 U.S. at 451--452}). Every Equal Protection case turns on three core inquiries: what class is harmed and whether it has faced a ``tradition of disfavor,'' what public purpose is served, and what characteristic of the class justifies disparate treatment. While many laws distinguishing persons with intellectual disabilities (such as special education or driving restrictions) would satisfy this impartial-lawmaker test, Cleburne's permit requirement---enforced amid neighborhood opposition and requiring signatures of nearby property owners---rested on irrational fears rather than any protective purpose a rational member of the disadvantaged class could approve (\href{https://supreme.justia.com/cases/federal/us/473/432/}{473 U.S. at 453--455}).

\subsection{Opinion Concurring in the Judgment in Part and Dissenting in Part (Justice Marshall)}

Justice Marshall, joined by Justices Brennan and Blackmun, agreed that the permit denial violated Equal Protection but sharply dissented from both the Court's methodology and its narrow as-applied remedy. First, Marshall criticized the majority's ``two for the price of one'' approach for violating constitutional avoidance by rejecting heightened scrutiny despite striking down the law on rational basis grounds. Second, he exposed the majority's unacknowledged departure from traditional deferential rational basis review (\textit{Williamson v. Lee Optical}), noting that under ``ordinary'' rationality review courts do not sift the factual record or invalidate underinclusive ``one step at a time'' legislation. Calling the Court's actual method ``second order'' rational basis review, Marshall warned that refusing to articulate the triggers for searching scrutiny risked both a return to \textit{Lochner}-era judicial activism in economic regulation and unprincipled, unpredictable lower-court decisions (\href{https://supreme.justia.com/cases/federal/us/473/432/}{473 U.S. at 456--460}).

Applying his sliding-scale Equal Protection framework (balancing the importance of the affected interest against the invidiousness of the classification), Marshall argued that heightened scrutiny was warranted because (1) establishing a community group home implicates a vital liberty interest, and (2) intellectually disabled persons suffered a ``lengthy and tragic history'' of state-mandated segregation, eugenic sterilization (\textit{Buck v. Bell}), marriage bans, school exclusion, and disenfranchisement that rivaled Jim Crow (\href{https://supreme.justia.com/cases/federal/us/473/432/}{473 U.S. at 460--465}). Marshall further argued that recent remedial legislation should reinforce rather than negate heightened scrutiny (as in \textit{Frontiero v. Richardson}) and that a trait's occasional relevance justifies intermediate rather than strict scrutiny. Finally, he dissented from the as-applied remedy, contending that leaving Cleburne's blunderbuss ``feeble-minded'' ordinance on the books departed from Equal Protection precedent and forced future applicants to run an unconstitutional gauntlet before biased local officials (\href{https://supreme.justia.com/cases/federal/us/473/432/}{473 U.S. at 465--478}).

\section{Examples: Future Applications}

\begin{enumerate}
  \item \textbf{Same-Side (Invalidated under \textit{Cleburne}):} A suburban municipality permits group residences for up to ten college students, religious order members, or recovering cardiac patients by right in multi-family zones, but requires an annual special use permit and neighbor consent for a supervised group home housing eight adults with autism. Because the city's asserted traffic, noise, and property-value concerns apply identically to freely permitted group dwellings, the differential burden fails \textit{Cleburne}'s rational basis review as resting on irrational prejudice.
  \item \textbf{Same-Side (Invalidated under \textit{Cleburne}):} A state statute bars any person with a history of depression or bipolar disorder from obtaining a cosmetology license, while permitting applicants with other chronic medical conditions to be licensed upon standard competency testing. Because the blanket exclusion bears no rational relationship to public health or safety and targets a disfavored group based on stereotypes rather than individualized capacity, it is invalid under \textit{Cleburne}.
  \item \textbf{Opposite-Side (Upheld under \textit{Cleburne}):} A state department of motor vehicles requires commercial bus drivers to pass an annual cognitive-reaction and medical fitness examination, resulting in the disqualification of applicants with moderate intellectual disabilities. Under \textit{Cleburne}, intellectual disability is not a quasi-suspect class, and ensuring highway safety is a legitimate state interest rationally related to objective cognitive testing.
  \item \textbf{Opposite-Side (Upheld under \textit{Cleburne}):} A city enforces a facially neutral single-family zoning ordinance limiting occupancy in R-1 zones to no more than four unrelated individuals, which incidentally prevents a nonprofit from opening an eight-resident group home for intellectually disabled adults. As the \textit{Cleburne} majority noted in Footnote~8 (citing \textit{Village of Belle Terre v. Boraas}), a neutral occupancy cap that treats disabled and non-disabled unrelated groups alike does not single out disability for irrational disfavor and survives rational basis review.
  \item \textbf{Fence-Sitter (Close Case):} A city requires a special use permit for any residential care facility---whether for intellectually disabled adults, recovering substance-use patients, or frail elderly residents---that houses more than eight people in a dense residential district, citing emergency-evacuation times for residents who require staff assistance during fires. Unlike Cleburne, the city treats comparable assisted-living populations alike and articulates a concrete safety rationale; yet challengers could argue that if unassisted dormitories or boarding houses of equal size are exempt without empirical proof of higher evacuation risk, the distinction still rests on unsubstantiated stereotypes rather than record evidence.
\end{enumerate}

\section{Critique}

\textit{Cleburne} is one of the most heavily analyzed Equal Protection decisions of the late twentieth century because of the glaring tension between its doctrinal label and its analytical rigor. Legal scholars famously dubbed the majority's approach ``rational basis with bite'' (\href{https://www.repository.law.indiana.edu/cgi/viewcontent.cgi?article=2074&context=ilj}{Gayle Lynn Pettinga, 62 Ind.\ L.J.\ 779 (1987)}; \href{https://www.scotusblog.com/2011/08/why-the-court-can-strike-down-marriage-restrictions-under-rational-basis-review/}{SCOTUSblog}). Traditional rational basis review defers to any conceivable legislative justification and allows legislatures to address problems ``one step at a time.'' Yet Justice White demanded record evidence from Cleburne, rejected conceivable safety and density rationales, and invalidated the ordinance for underinclusiveness. This doctrinal sleight of hand left lower courts without clear criteria for determining when ordinary deferential rationality ends and searching ``second order'' review begins, a tension that later resurfaced in animus cases such as \textit{Romer v. Evans} and \textit{United States v. Windsor}.

From a \textbf{progressive and antisubordination perspective}, critics echo Justice Marshall's charge that the Court minimized a brutal history of eugenics, forced institutionalization, and civic exclusion. By treating recent remedial statutes as proof that people with intellectual disabilities are not ``politically powerless,'' the Court created a perverse paradox: civil rights progress in legislatures strips a historically marginalized group of heightened constitutional protection, a logic never applied to race or sex. Moreover, \textit{Cleburne}'s refusal to recognize heightened scrutiny had major downstream consequences when the Supreme Court later relied on \textit{Cleburne} in \textit{Board of Trustees of the University of Alabama v. Garrett} (2001) to curtail Congress's Section~5 power to abrogate state sovereign immunity under Title~I of the Americans with Disabilities Act (\href{https://journals.law.harvard.edu/jol/2025/02/22/our-rhizomatic-constitution/}{Harvard Journal on Legislation}). Critics also fault the as-applied remedy for leaving the stigmatizing ``feeble-minded'' text intact.

From an \textbf{originalist and textualist perspective}, \textit{Cleburne} draws two competing lines of critique. Traditional judicial-restraint conservatives and formalists argue that if rational basis review is truly the governing standard, the majority overstepped its institutional role by second-guessing local zoning judgments and inferring ``irrational prejudice'' from legislative underinclusiveness, thereby resurrecting open-ended judicial policymaking without textual mooring. Conversely, classical-liberal and libertarian originalists who question the constitutional legitimacy of hyper-deferential rational basis review contend that the Fourteenth Amendment's text and the Article~III judicial duty require courts to evaluate actual statutory ends and real-world evidence rather than inventing hypothetical rationales for the government (\href{https://digitalcommons.law.villanova.edu/cgi/viewcontent.cgi?article=3486&context=vlr}{Joseph S. Diedrich, 66 Vill.\ L. Rev.\ 249 (2021)}). From that vantage point, \textit{Cleburne}'s searching factual inquiry was substantively sound, but the Court erred by disguising it as an ad hoc exception rather than candidly restoring meaningful judicial review of arbitrary legislation.

\section{Key Quotations}

\begin{quote}
The Equal Protection Clause of the Fourteenth Amendment commands that no State shall ``deny to any person within its jurisdiction the equal protection of the laws,'' which is essentially a direction that all persons similarly situated should be treated alike.
\end{quote}
\noindent --- Justice White, Majority Opinion, \href{https://supreme.justia.com/cases/federal/us/473/432/}{473 U.S. at 439}.

\begin{quote}
But mere negative attitudes, or fear, unsubstantiated by factors which are properly cognizable in a zoning proceeding, are not permissible bases for treating a home for the mentally retarded differently from apartment houses, multiple dwellings, and the like.
\end{quote}
\noindent --- Justice White, Majority Opinion, \href{https://supreme.justia.com/cases/federal/us/473/432/}{473 U.S. at 448}.

\begin{quote}
The short of it is that requiring the permit in this case appears to us to rest on an irrational prejudice against the mentally retarded, including those who would occupy the Featherston facility and who would live under the closely supervised and highly regulated conditions expressly provided for by state and federal law.
\end{quote}
\noindent --- Justice White, Majority Opinion, \href{https://supreme.justia.com/cases/federal/us/473/432/}{473 U.S. at 450}.

\begin{quote}
Rather, our cases reflect a continuum of judgmental responses to differing classifications which have been explained in opinions by terms ranging from ``strict scrutiny'' at one extreme to ``rational basis'' at the other.
\end{quote}
\noindent --- Justice Stevens, Concurring Opinion, \href{https://supreme.justia.com/cases/federal/us/473/432/}{473 U.S. at 451}.

\begin{quote}
To be sure, the Court does not label its handiwork heightened scrutiny, and perhaps the method employed must hereafter be called ``second order'' rational basis review, rather than ``heightened scrutiny.'' But however labeled, the rational basis test invoked today is most assuredly not the rational basis test of \textit{Williamson v. Lee Optical of Oklahoma, Inc.}, 348 U.S. 483 (1955); \textit{Allied Stores of Ohio, Inc. v. Bowers}, 358 U.S. 522 (1959), and their progeny.
\end{quote}
\noindent --- Justice Marshall, Concurring in the Judgment in Part and Dissenting in Part, \href{https://supreme.justia.com/cases/federal/us/473/432/}{473 U.S. at 458}.


\end{document}
